% Stage D2 state-ingestion report snippet.
\documentclass[11pt,a4paper]{article}
\usepackage[margin=25mm]{geometry}
\usepackage{booktabs}
\usepackage{hyperref}
\title{Stage D2: Minimal Abaqus State-Ingestion Verification}
\author{Adaptive Remeshing Thesis}
\date{2026-07-22}
\begin{document}
\maketitle

\section{Status}
The D2A executable package and HPC ingestion job passed with classification
\path{stage_d2a_state_ingestion_pass}. The accepted PBS job was
\texttt{1376785.mmaster02}; it finished with PBS \texttt{Exit\_status=0},
solver exit zero, readable ODB output, and a written \texttt{D2A.ok} marker.

\section{D1 transfer-error baselines}
\begin{table}[h]
\centering
\begin{tabular}{@{}lr@{}}
\toprule
Quantity & Value \\
\midrule
Nodal $d$ L2 error & $0.0270$ \\
Nodal $d$ maximum error & $0.0850$ \\
IP $H$ L2 error & $0.0108$ \\
IP $H$ maximum error & $0.0234$ \\
Energy difference & $-0.00769$ \\
\bottomrule
\end{tabular}
\end{table}

\section{D2A route}
The D2A source variant confirms the preserved Molnar U1 phase DOF as DOF~3 and
uses a generated element/IP keyed table for transferred history initialization.
The tiny target contains $15$ target nodes, $8$ physical elements, and one
visualization integration point per physical element. The phase nodes use
labels $1$--$15$, while the mechanics and visualization nodes are duplicated as
labels $1001$--$1015$. The U1 layer uses element labels $1$--$8$, the U2 layer
uses $9$--$16$, and the visualization layer uses $17$--$24$.

\section{Accepted result}
\begin{table}[h]
\centering
\begin{tabular}{@{}lr@{}}
\toprule
Quantity & Value \\
\midrule
Job ID & \texttt{1376785.mmaster02} \\
PBS exit status & $0$ \\
Target element/IP coverage & $1.0$ \\
Maximum \texttt{SDV15} interpolation error & $0.0$ \\
Maximum \texttt{SDV16}/$H$ error & $6.428999999030793\times10^{-9}$ \\
Failures & $0$ \\
\bottomrule
\end{tabular}
\end{table}

The ODB comparison verifies transferred phase ingestion through the UMAT mirror
\texttt{SDV15}, which equals the independent interpolation of transferred nodal
phase on every target element/IP. It verifies transferred history ingestion
through \texttt{SDV16}, which matches transferred $H$ within the $10^{-8}$
tolerance. Abaqus did not expose the UEL phase DOF as usable nodal
\texttt{U} output in this smoke deck, so nodal \texttt{U} is recorded as an ODB
output limitation rather than as the accepted phase-ingestion signal.

\section{D2B Serial Continuation Attempt}
The next serial continuation gate was submitted exactly once as PBS job
\texttt{1376819.mmaster02}. It did not pass. Abaqus compiled the user
subroutine and entered the three-step continuation deck, but the job finished
with PBS \texttt{Exit\_status=10} and solver exit $1$. The message file reports
\texttt{TOO MANY INCREMENTS NEEDED TO COMPLETE THE STEP} during the tiny
mechanical continuation. No \texttt{D2B.ok} marker was written.

\begin{table}[h]
\centering
\begin{tabular}{@{}lr@{}}
\toprule
Quantity & Value \\
\midrule
Job ID & \texttt{1376819.mmaster02} \\
PBS exit status & $10$ \\
Solver exit & $1$ \\
Classification & \path{stage_d2b_solver_fail} \\
Accepted marker & no \texttt{D2B.ok} \\
\bottomrule
\end{tabular}
\end{table}

A corrected D2B deck has been prepared for a future authorized rerun by
increasing the release and continuation step increment allowance to
\texttt{inc=50}. This correction does not change transferred field values,
transfer-table values, label mapping, or the D2-only UEL/UMAT state route.

\end{document}
